\textbf{Phần 1. Trắc nghiệm lựa chọn}
\setcounterpageref{ex}{0}

\begin{ex} %[3D2V40]
Tính giá trị của biểu thức: $A = (150 + 250) : 5 - 35$.
\choice
{$40$}
{\True $45$}
{$50$}
{$55$}
\loigiai{
$400 : 5 - 35 = 80 - 35 = 45$.
\textbf{Đáp án của câu hỏi là B}
}
\end{ex}

\begin{ex} %[3D2V40]
Một cửa hàng có 3 thùng táo, mỗi thùng có 120 quả táo. Người ta đã bán được $\dfrac{1}{4}$ tổng số táo. Cửa hàng còn lại bao nhiêu quả táo?
\choice
{$90$ quả}
{$250$ quả}
{$280$ quả}
{\True $270$ quả}
\loigiai{
Tổng số táo: $120 \times 3 = 360$ quả. Đã bán: $360 : 4 = 90$ quả. Còn lại: $360 - 90 = 270$ quả.
\textbf{Đáp án của câu hỏi là D}
}
\end{ex}

\begin{ex} %[3D2V40]
Tìm $x$, biết: $x \times 6 - 45 = 555$.
\choice
{\True $100$}
{$95$}
{$105$}
{$90$}
\loigiai{
$x \times 6 = 555 + 45 = 600 \Rightarrow x = 600 : 6 = 100$.
\textbf{Đáp án của câu hỏi là A}
}
\end{ex}

\begin{ex} %[3D2V40]
Số thứ nhất là 180, số thứ hai bằng $\dfrac{1}{3}$ số thứ nhất. Hỏi số thứ nhất nhiều hơn số thứ hai bao nhiêu đơn vị?
\choice
{$60$}
{$90$}
{\True $120$}
{$150$}
\loigiai{
Số thứ hai là: $180 : 3 = 60$. Nhiều hơn: $180 - 60 = 120$ đơn vị.
\textbf{Đáp án của câu hỏi là C}
}
\end{ex}

\begin{ex} %[3D2V40]
Một cuộn dây thép dài $450\text{ m}$. Người ta cắt đi 3 đoạn, mỗi đoạn dài $110\text{ m}$. Cuộn dây còn lại dài bao nhiêu mét?
\choice
{$110\text{ m}$}
{\True $120\text{ m}$}
{$130\text{ m}$}
{$150\text{ m}$}
\loigiai{
Đã cắt: $110 \times 3 = 330\text{ m}$. Còn lại: $450 - 330 = 120\text{ (m)}$.
\textbf{Đáp án của câu hỏi là B}
}
\end{ex}

\begin{ex} %[3D2V40]
Một hình chữ nhật có chu vi $64\text{ cm}$, chiều rộng là $8\text{ cm}$. Chiều dài gấp mấy lần chiều rộng?
\choice
{$2$ lần}
{$4$ lần}
{$5$ lần}
{\True $3$ lần}
\loigiai{
Nửa chu vi: $64 : 2 = 32\text{ cm}$. Chiều dài: $32 - 8 = 24\text{ cm}$. Chiều dài gấp chiều rộng: $24 : 8 = 3$ (lần).
\textbf{Đáp án của câu hỏi là D}
}
\end{ex}

\begin{ex} %[3D2V40]
Trong phép chia cho 7, nếu tăng số bị chia thêm 14 đơn vị và giữ nguyên số chia thì thương sẽ
\choice
{\True Tăng thêm 2 đơn vị}
{Tăng thêm 14 đơn vị}
{Tăng thêm 7 đơn vị}
{Không đổi}
\loigiai{
Vì $14 : 7 = 2$ nên thương sẽ tăng thêm đúng 2 đơn vị.
\textbf{Đáp án của câu hỏi là A}
}
\end{ex}

\begin{ex} %[3D2V40]
Tính giá trị biểu thức: $750 : 5 + 125 \times 4$.
\choice
{$600$}
{$700$}
{\True $650$}
{$550$}
\loigiai{
$750 : 5 = 150$; $125 \times 4 = 500$; $150 + 500 = 650$.
\textbf{Đáp án của câu hỏi là C}
}
\end{ex}

\begin{ex} %[3D2V40]
Bác nông dân nuôi 5 chuồng gà, mỗi chuồng có 85 con gà. Sau khi bán đi một số con gà thì số gà còn lại bằng $\dfrac{1}{5}$ tổng số gà ban đầu. Bác nông dân đã bán bao nhiêu con gà?
\choice
{$85$ con gà}
{\True $340$ con gà}
{$425$ con gà}
{$320$ con gà}
\loigiai{
Tổng số gà: $85 \times 5 = 425$ con. Số gà còn lại: $425 : 5 = 85$ con. Số gà đã bán: $425 - 85 = 340$ con.
\textbf{Đáp án của câu hỏi là B}
}
\end{ex}

\begin{ex} %[3D2V40]
Tìm số bị chia trong phép chia có số chia là 9, thương là 86 và số dư lớn nhất có thể.
\choice
{$774$}
{$783$}
{$781$}
{\True $782$}
\loigiai{
Số dư lớn nhất là 8. Số bị chia: $9 \times 86 + 8 = 774 + 8 = 782$.
\textbf{Đáp án của câu hỏi là D}
}
\end{ex}

\begin{ex} %[3D2V40]
Một mảnh vườn hình chữ nhật có chiều dài $120\text{ m}$, chiều rộng bằng $\dfrac{1}{3}$ chiều dài. Chu vi của mảnh vườn đó là
\choice
{\True $320\text{ m}$}
{$240\text{ m}$}
{$360\text{ m}$}
{$280\text{ m}$}
\loigiai{
Chiều rộng: $120 : 3 = 40\text{ m}$. Chu vi: $(120 + 40) \times 2 = 320\text{ (m)}$.
\textbf{Đáp án của câu hỏi là A}
}
\end{ex}

\begin{ex} %[3D2V40]
Tổng của hai số là 150, số lớn gấp 4 lần số bé. Số lớn là
\choice
{$30$}
{$100$}
{\True $120$}
{$110$}
\loigiai{
Tổng số phần bằng nhau: $4 + 1 = 5$ phần. Số bé: $150 : 5 = 30$. Số lớn: $30 \times 4 = 120$.
\textbf{Đáp án của câu hỏi là C}
}
\end{ex}

\textbf{Phần 2. Trắc nghiệm đúng sai}
\setcounterpageref{ex}{0}

\begin{ex} %[3D2V40]
Xét tính đúng sai của các phép toán tổng hợp:
\choiceTF
{\True Biểu thức $(100 - 20) : 4$ có giá trị bằng $20$.}
{Biểu thức $100 - 20 : 4$ có giá trị bằng $20$.}
{\True Số 100 gấp 5 lần số 20.}
{Số 100 nhiều hơn số 20 là 5 đơn vị.}
\loigiai{
\begin{itemize}
    \item Ý a) Đúng: $80 : 4 = 20$.
    \item Ý b) Sai: Sai vì $100 - 5 = 95$.
    \item Ý c) Đúng: $100 : 20 = 5$.
    \item Ý d) Sai: Sai vì $100 - 20 = 80$ đơn vị.
\end{itemize}
\begin{itemize}
    \item a) \textbf{Đúng}.
    \item b) \textbf{Sai}.
    \item c) \textbf{Đúng}.
    \item d) \textbf{Sai}.
\end{itemize}
}
\end{ex}

\begin{ex} %[3D2V40]
Một cửa hàng có 3 kho gạo, mỗi kho chứa 250 bao gạo:
\choiceTF
{\True Tổng số gạo trong cả ba kho là 750 bao.}
{\True Nếu bán đi 1 kho thì còn lại 500 bao gạo.}
{Số gạo còn lại sau khi bán 1 kho gấp 3 lần số gạo đã bán.}
{\True Mỗi bao gạo nặng $50\text{ kg}$ thì cả 3 kho chứa tổng cộng $37500\text{ kg}$ gạo.}
\loigiai{
\begin{itemize}
    \item Ý a) Đúng: $250 \times 3 = 750$ bao.
    \item Ý b) Đúng: $750 - 250 = 500$ bao.
    \item Ý c) Sai: Sai vì 500 chỉ gấp đôi 250.
    \item Ý d) Đúng: Phép nhân khối lượng.
\end{itemize}
\begin{itemize}
    \item a) \textbf{Đúng}.
    \item b) \textbf{Đúng}.
    \item c) \textbf{Sai}.
    \item d) \textbf{Đúng}.
\end{itemize}
}
\end{ex}

\textbf{Phần 3. Trả lời ngắn}
\setcounterpageref{ex}{0}

\begin{ex} %[3D2V40]
Tính giá trị của biểu thức: $600 : 3 + 120 \times 2$.
\shortans{440}
\loigiai{
\textbf{Đáp số:} $440$\\
$200 + 240 = 440$.
}
\end{ex}

\begin{ex} %[3D2V40]
Một đoàn gồm 4 xe tải, mỗi xe chở 125 bao xi măng. Người ta chuyển đi 220 bao xi măng. Hỏi còn lại bao nhiêu bao xi măng?
\shortans{280}
\loigiai{
\textbf{Đáp số:} $280$\\
Tổng số bao: $125 \times 4 = 500$ bao. Còn lại: $500 - 220 = 280$ bao.
}
\end{ex}

\begin{ex} %[3D2V40]
Số lớn là 96, số bé là 8. Hỏi số lớn gấp mấy lần số bé?
\shortans{12}
\loigiai{
\textbf{Đáp số:} $12$\\
$96 : 8 = 12$ lần.
}
\end{ex}

\begin{ex} %[3D2V40]
Tìm số tự nhiên $y$, biết: $(y - 50) : 4 = 100$.
\shortans{450}
\loigiai{
\textbf{Đáp số:} $450$\\
$y - 50 = 400 \Rightarrow y = 450$.
}
\end{ex}

\begin{ex} %[3D2V40]
Một cuộn dây dài $720\text{ m}$. Người ta cắt đi $\dfrac{1}{8}$ cuộn dây. Hỏi cuộn dây còn lại dài bao nhiêu mét?
\shortans{630}
\loigiai{
\textbf{Đáp số:} $630$\\
Đã cắt: $720 : 8 = 90\text{ m}$. Còn lại: $720 - 90 = 630\text{ (m)}$.
}
\end{ex}

\begin{ex} %[3D2V40]
Chu vi của một hình vuông bằng chu vi hình tam giác đều có cạnh $16\text{ cm}$. Tính độ dài cạnh của hình vuông đó theo xăng-ti-mét.
\shortans{12}
\loigiai{
\textbf{Đáp số:} $12$\\
Chu vi tam giác: $16 \times 3 = 48\text{ cm}$. Cạnh hình vuông: $48 : 4 = 12\text{ (cm)}$.
}
\end{ex}
